\chapter{Matriz de identidad profesional}

Esta parte de la asignatura se centró en el autoconocimiento como punto de partida para tomar decisiones coherentes sobre la propia carrera profesional. La idea que vertebró toda la sesión fue que, antes de pensar en sectores, empresas o trayectorias, es necesario entender quiénes somos, qué queremos y qué necesitamos para desarrollarnos de forma sostenible a lo largo del tiempo.

La profesora comenzó recordando que las emociones están presentes en todas nuestras decisiones, también en las profesionales. No actuamos de manera fría o perfectamente racional, por lo que aprender a reconocer cómo influyen nuestras emociones, motivaciones y patrones de comportamiento es esencial para construir un proyecto profesional estable. A partir de esta idea, introdujo el concepto de competencias emocionales, que ayudan precisamente a identificar qué nos impulsa, qué nos bloquea y qué condiciones necesitamos para rendir bien.

A continuación trabajamos con la distinción entre sueño, visión y misión. El sueño es aquello que deseamos pero que aún no hemos transformado en un objetivo concreto. La visión, en cambio, implica asumir el coste, hacerla propia y plantearla como algo alcanzable. La misión supone traducir esa visión en acciones reales. Esta diferencia resultó útil para ordenar ideas y comprender que avanzar profesionalmente no depende solo de tener metas, sino de convertirlas en decisiones prácticas.

En este contexto, se presentaron varias herramientas de planificación personal. Una de ellas fue el método de retroplanificación, que consiste en partir del objetivo final e ir retrocediendo en el tiempo para identificar los pasos previos necesarios. Esta técnica ayuda a dividir metas ambiciosas en acciones pequeñas y manejables. También se explicó la matriz de Eisenhower, que distingue entre lo importante y lo urgente. La reflexión principal fue que, para desarrollarse profesionalmente, es necesario reservar tiempo a tareas importantes pero no urgentes, como formarse, reflexionar, explorar sectores o planificar la trayectoria.

Con este marco, pasamos a trabajar la matriz de identidad profesional, pensada para ordenar cuatro dimensiones fundamentales: el tipo de organización en la que nos gustaría trabajar, el sector de actividad, el ámbito funcional y la orientación vocacional. El objetivo no era elegir una única opción, sino identificar aquellas que realmente encajan con nuestra personalidad, intereses y forma de trabajar. Fue interesante observar cómo preferencias aparentemente personales —como disfrutar del trabajo con personas o buscar proyectos con impacto social— pueden convertirse en criterios relevantes a la hora de orientar la carrera.

En los apartados de fortalezas y áreas de mejora se recordó que ambos elementos deben considerarse de manera honesta. No se trata de elaborar una lista idealizada, sino de reconocer tanto aquello que ya sabemos hacer bien como lo que aún necesitamos mejorar para acercarnos a nuestros objetivos. Este ejercicio de honestidad es clave para evitar decisiones automáticas o basadas en expectativas externas.

El último componente de la matriz es el plan, donde cada persona concreta los pasos que debería dar para avanzar hacia su visión profesional. Puede incluir acciones formativas, experiencias que adquirir, competencias que reforzar o decisiones que postergar. La utilidad de la matriz no reside en cerrarla de forma definitiva, sino en revisarla periódicamente para adaptarla a los cambios personales o del entorno.

En conjunto, esta dinámica resultó muy valiosa porque ofreció un método claro para aterrizar ideas que normalmente quedan dispersas. Comprendí que construir una trayectoria profesional no depende solo de elegir un sector o un puesto, sino de conocerse, priorizar y tomar decisiones con intención. Tener una visión clara facilita enfocar la energía, evitar distracciones y avanzar hacia un camino alineado con lo que realmente buscamos a medio y largo plazo.
